\subsection{Buying the full-mesh finite-difference bounds}
\label{sec:ex-fullmesh-fixedbounds}

We prove the finite-volume estimate needed before any continuum extraction. The
important order is to pay for the bounded projected coefficient first, and only
then use the exact Wilson difference identity. The result is uniform over every
period at all sufficiently fine meshes. The proof has two parts: a finite
dimensional absorption estimate for one coefficient row, and a packing argument
which plants finitely many such rows into one large periodic box.

For a period $p$, mesh $N$, field list $v=(v_1,\ldots,v_n)$, and finite
list $l$ of slot-and-coordinate directions, write $K_{N,p;v,l}(x)$ for the
finite Wilson expectation of the product of the corresponding iterated
periodic difference insertions at the sites $x=(x_1,\ldots,x_n)$ of the
period-$p$ torus at mesh $N$. Sites have integer coordinates; the physical
position of a site $z$ is $a_Nz$. The normalization is
$\Delta_{N,\alpha}F(z)=a_N^{-1}\bigl(F(z+e_\alpha)-F(z)\bigr)$ for a site $z$
and a unit lattice step $e_\alpha$. Thus the theorem below concerns the
unscaled derivative kernel; the exact one-step identity used in its proof is
$F(z+e_\alpha)-F(z)=a_N\Delta_{N,\alpha}F(z)$.

Throughout this section the index set $I$ is the set of mesh-period pairs
$i=(N,p)$ with the preorder $(N,p)\le(N',p')$ if and only if $N\le N'$.
A statement holds \emph{eventually in $I$} if there is $N_0$ such that it
holds for all $N\ge N_0$ and every period $p$. A limit in $I$ is understood
in the same sense: $y_i\to y$ means that for every $\varepsilon>0$ there is
$N_0$ with $|y_{(N,p)}-y|<\varepsilon$ for all $N\ge N_0$ and all $p$.
With the product order instead, ``eventually'' would exclude small periods
and the uniformity asserted below would not follow.

\begin{lemma}[Absorption of the projected coefficient]\label{lem:ex-fullmesh-cap}
Let $I$ be the mesh-period index set just described, and let $X_i$ be a
finite-dimensional normed coefficient space. Suppose that for each $i$ there
are a projection $P_i:X_i\to X_i$, a linear map $L_i:X_i\to Y_i$, and a constant $Z>0$
such that
\begin{equation}
   \|P_i z\|\le Z\|L_i(P_i z)\|.                         \label{eq:ex-cap-rank}
\end{equation}
Let $x_i\in X_i$, put $X_i^0=P_i x_i$, and set
$\lambda_i=(1+\|X_i^0\|)^{-1}$. Assume the following estimates hold
eventually in $I$:
\begin{align}
 \|L_i x_i\|&\le c,                                         \label{eq:ex-cap-tested}\\
 \|L_i(\lambda_i(x_i-X_i^0))\|&\le \frac{1}{2Z}.             \label{eq:ex-cap-remainder}
\end{align}
Then eventually
\begin{equation}
                    \|X_i^0\|\le 2Zc+1.                    \label{eq:ex-cap-conclusion}
\end{equation}
\end{lemma}

\begin{proof}
Write $R_i=x_i-X_i^0$. By linearity and \eqref{eq:ex-cap-tested},
\[
 \|X_i^0\|\le Z\|L_iX_i^0\|
 \le Zc+Z\|L_iR_i\|.
\]
Since $\lambda_i=(1+\|X_i^0\|)^{-1}$, \eqref{eq:ex-cap-remainder} implies
$Z\|L_iR_i\|\le(1+\|X_i^0\|)/2$. Therefore
\[
 \|X_i^0\|\le Zc+\tfrac12(1+\|X_i^0\|),
\]
and rearranging gives \eqref{eq:ex-cap-conclusion}. All estimates were eventual, so their two mesh
thresholds can be replaced by their maximum.
\end{proof}

\begin{proposition}[One coefficient row is bounded at every period]
\label{prop:ex-fullmesh-cap-source}
Let $C\subseteq\prod_i X_i$ be one class of coefficient row families satisfying the following normalized-cost conditions. Suppose that the full family
$x=(x_i)$ belongs to $C$, that $C$ is closed under subtraction and
pointwise multiplication by scalar families whose absolute values are at most
one, and that every uniformly norm-bounded row family belongs to $C$. Suppose also that the detector estimates
\begin{enumerate}
\item $\|P_i z\|\le Z\|L_i(P_i z)\|$ for all $z\in X_i$ with a
      scale-independent $Z>0$;
\item for the primary paid trial there are assay values $b_i\in Y_i$ with
      $\|b_i\|\le v$ and $\|L_i x_i-b_i\|\to0$;
\item for every row family $z\in C$ satisfying $P_i z_i=0$ for every $i$,
      one has $\|L_i z_i\|\to0$.
\end{enumerate}
Then, eventually in $I$, that is, for all sufficiently fine meshes and
uniformly over all periods,
\[
                        \|P_i x_i\|\le 2Z(v+1)+1.
\]
The uniformity means that the limits in items~2 and~3 are taken in $I$:
they hold beyond one mesh threshold for every period.
\end{proposition}

\begin{proof}
From item~2, eventually in $I$, $\|L_i x_i\|\le v+1=:c$. Put
$X_i^0=P_i x_i$ and
\[
       \lambda_i=(1+\|X_i^0\|)^{-1},\qquad
       z_i=\lambda_i(x_i-X_i^0).
\]
The scalar family $\lambda=(\lambda_i)_i$ satisfies $|\lambda_i|\le1$, so
$\lambda x\in C$. The family $\lambda X^0$ is uniformly bounded by one, so
it belongs to $C$ as well. Therefore $z=\lambda(x-X^0)\in C$ by subtraction,
and $P_i z_i=0$ because
$P_i(x_i-P_ix_i)=0$. By item~3,
for all sufficiently fine meshes $\|L_i z_i\|\le(2Z)^{-1}$. Since
$L_i z_i=\lambda_i L_i(x_i-X_i^0)$, estimate \eqref{eq:ex-cap-rank} gives
\[
 \|X_i^0\|\le Z\|L_iX_i^0\|
 \le Zc+Z\|L_i(x_i-X_i^0)\|
 \le Zc+\frac{1+\|X_i^0\|}{2}.
\]
Rearranging proves the claim. The common mesh threshold is the maximum of the
mesh thresholds in items~2 and~3.
\end{proof}

The cap alone bounds a projected row. The difference trial is paid by a
separate closure argument which uses that bound; we give its steps. The
primary paid trial, the exact equation with its vanishing owed cost, the
bounded-coefficient receipt and the sweep for rows with zero projection are
first kept explicit in
Theorem~\ref{thm:ex-fullmesh-fixedbounds-producer} below.
Proposition~\ref{prop:fullmesh-selected-supplier} verifies them for the
retained public family. Let $x$ be
the coefficient row for the requested insertion and let $F$ be its primary
paid trial. The exact equation gives an owed row
$e=\operatorname{upper}(x)-\operatorname{hops}(F)$, with a root-test cost
$u$ tending to zero. Applying the cap to $(x,F,e)$ bounds $Px$ uniformly.
Thus $y=Px$ is a uniformly bounded row family and its primary insertion is
paid by the bounded-coefficient receipt. The row $x-y$ lies in the normalized
cost class by the class closure, and $P(x-y)=0$. The sweep theorem therefore
makes $\operatorname{upper}(x-y)$ owed with a cost tending to zero. We
have the identity
\[
  \operatorname{hops}(F)
   =\operatorname{upper}(y)+\operatorname{upper}(x-y)-e.
\]
It follows from $\operatorname{upper}(x)=\operatorname{upper}(y)+
\operatorname{upper}(x-y)$ and the defining equation for $e$. The paid
budget is closed under finite addition and subtraction: once each vanishing
cost is below the unit allowance at a sufficiently fine mesh, the primary
paid budget together with the two unit owed budgets yields a finite paid
budget for the requested difference trial. This is exactly the projection-and-remainder argument;
it does not require the unsplit equation residual to be bounded by a quantity
that depends on $\|Px\|$.

\begin{theorem}[Full-mesh bounds for all finite differences]
\label{thm:ex-fullmesh-fixedbounds-producer}
Fix a family for which the preceding coefficient-row estimate and the
paid-remainder decomposition just described apply to every finite list of
fields and difference directions, including each resulting difference trial
transformed by a signed-coordinate action, with or without coefficientwise
conjugation. Assume that each bounded primary coefficient row family and each
of its finitely many corresponding orientations has a paid root-test
budget in the following sense: for each such oriented row family there is a
depth $k_0$ such that, for every centred scale $k\ge k_0$, there are a
budget $b<\infty$ and a mesh threshold valid for every period beyond which
the centred root test at scale $k$ is at most $b$. The bounds and thresholds
may depend on the oriented family and $k$. Assume also that the exact
periodic one-step difference identity holds, and that translated insertions satisfy the product root-test
inequality in distinct interior cells.
Then for each finite field list $v$, difference list $l$, and $r>0$, there
are $B<\infty$ and a mesh threshold $N_0$ such that for every $N\ge N_0$,
every period $p$, and every periodic position tuple $x$ such that any two
distinct slot indices $u\ne u'$ have sites $x_u,x_{u'}$ with
\[
 a_N\,\operatorname{sep}_{N,p}(x_u,x_{u'})\ge r
\]
(even if the corresponding insertions agree),
\begin{equation}
                         |K_{N,p;v,l}(x)|\le B.              \label{eq:ex-fullmesh-bound}
\end{equation}
The threshold is independent of $p$; the bound is independent of both $N$
and $p$. Here $\operatorname{sep}_{N,p}(y,y')=\max_{\alpha}|y_\alpha-y'_\alpha|_{\rm per}$
is the periodic maximum-coordinate distance of two sites of the period-$p$
torus at mesh $N$, where $|k|_{\rm per}$ denotes the distance from the integer
$k$ to the nearest multiple of the side length $\ell_{N,p}$. Thus the
separation is measured in the periodic $\ell^\infty$ metric, not the
Euclidean one; the injectivity argument below uses exactly this metric.
\end{theorem}

\begin{proof}
First fund one row. For a field label $w$, the constant coefficient row
$x_i=w$ has bounded normalized cost. The assumed primary receipt and the
preceding cap-and-remainder argument give a paid budget for each finite
difference list applied to that insertion, eventually at every mesh and
every period.
For a general finite list $v=(v_1,\ldots,v_n)$ and positional derivative
list, make one trial for each slot by applying to $v_j$ the subsequence of
directions assigned to that slot. The exact one-step identity is the
periodic discrete difference formula
\[
 F_j(z+e_\alpha)-F_j(z)
   =a_N\,\Delta_{N,\alpha}F_j(z),
\]
with the normalization of the difference operator. Iterating this identity
and the exact product measure definition gives the desired multi-difference
kernel, including wrapped coordinates. The list is finite, so the paid
records have finitely many depth thresholds $k_0$.

We now explain the two-scale tiling argument. First choose a tiling scale
$\ell$ with $a_\ell<r$. Then use $a_k\to0$ to choose one recentering scale
$k\ge\ell$ which lies beyond all the finitely many depth thresholds $k_0$ of
the paid records and satisfies
\[
                         6n a_k<a_\ell .
\]
At this scale $k$ all the root-test budgets hold beyond a common mesh
threshold, uniformly over periods (take the maximum of the finitely many
mesh thresholds). At a fine mesh $N$ beyond both scales, let $t$ be the
integer span of scale $k$ and let $T$ be a periodic tiling whose side is the
integer span of scale $\ell$. The aligned-scale identities give
$a_Nt=a_k$ and $a_N T.\mathrm{side}=a_\ell$. Thus
$n(2(3t))<T.\mathrm{side}$. The tiling placement lemma, applied with margin
$3t$, gives one translated tiling and an assignment of the $n$ insertion
sites to cells, each site lying at least $3t$ from the cell boundary.
The boundary count is the one in the proof of
Theorem~\ref{ex-centered-producers}: in each coordinate one site forbids at
most $6t$ of the $T.\mathrm{side}$ possible translations, so the $n$ sites
forbid fewer than $T.\mathrm{side}$ of them.
This is uniform in the period: alignment gives the same physical side
$a_\ell$ for every period index, once $N$ is sufficiently fine.

The cell assignment is injective for a geometric reason. If two sites were
assigned to the same cell, then in each coordinate both lie in one window of
$T.\mathrm{side}$ consecutive residues, so their periodic coordinate
distance is at most $T.\mathrm{side}$; hence
$\operatorname{sep}_{N,p}\le T.\mathrm{side}$.
Multiplying by $a_N$ would give $a_N\operatorname{sep}_{N,p}\le
a_NT.\mathrm{side}=a_\ell<r$, contradicting the
assumption. (With Euclidean separation one would instead need
$2a_\ell<r$ in four dimensions; the statement above avoids this by using the
$\ell^\infty$ metric throughout.) The three-span margin is exactly what the centering lemma needs
to transport each local insertion from the centered tile to its assigned
cell.

For each occupied cell, recenter its insertion to the distinguished center
cell. The budgets are root tests relative to the fine centred tiling at
scale $k$, whereas the planted inequality below is applied in the coarse
tiling of side $a_\ell$. Invariance of the Wilson law under translations and
signed coordinate symmetries does not by itself compare root tests of
different tilings. The comparison is the replicated-orbit and
maximum-principle argument in the proof of
Theorem~\ref{ex-centered-producers} (the packing argument, with its coarse
depth $m$ and fine depth $k$ taken to be $\ell$ and $k$ here): the fine
planted inequality bounds the root test of each insertion relative to the
coarse tiling at a home point on the fine centre grid by the fine budget,
the one-dimensional maximum principle moves the home point to the assigned
site, and signed-coordinate covariance identifies the oriented copy in each
occupied cell with the original insertion. The hypothesis supplies the
root-test budgets for the finitely many signed-coordinate and conjugated
orientations of each trial. Include their paid records in the finite list
above, as in Lemma~\ref{ex-centered-finite}, before taking the maximum of the
depth thresholds and the maximum of the mesh thresholds at the chosen scale.
This transports both the
interior support condition and the same root-test budget to the occupied
cell. No statistical independence is asserted: all factors
are evaluated under the same Wilson law. Instead, the planted-cell
inequality applies to the translated copies in distinct cells and gives
\[
 \left|\mathbb E\prod_{j=1}^n F_j(x_j)\right|
       \leq \prod_{j=1}^n b_j,
\]
where $b_j$ is the paid root-test bound for the $j$th trial. Their finite
product is a constant independent of mesh and period. This proves
\eqref{eq:ex-fullmesh-bound}.
\end{proof}

\subsubsection{The uniform supplier on the retained family}

We now verify the hypotheses above on $C_R$. In this subsection every
occurrence of ``eventually'' means eventually in $I$, including periods
smaller than any prescribed comparison level. This is stronger than an
ordered extraction statement.

Call a row $x=(x_{N,p})$ of grade $d$ \emph{uniformly subpower} if
\[
 \forall\epsilon>0\ \exists N_\epsilon\ \forall N\ge N_\epsilon\
 \forall p,\qquad |X_{N,p}^{-1}x_{N,p}|_1\le e^{\epsilon s_N}.
\]
Here $X$ is the actual terminal calibration of
Subsection~\ref{sec:actual-equation-kernel-relay}. Denote this class by
$\mathscr C_d$. It is closed under addition, subtraction and scalar
multiplication by families of modulus at most one: split the target
exponent between the finitely many prices and absorb their fixed sum
using $s_N\to\infty$. Every norm-bounded row belongs to it. Indeed the
reverse-column estimate of Subsubsection~\ref{ex-column-subpower} bounds its
inverse coordinates by a constant times $J_N$, and the displayed
proof $\log J_N=o(s_N)$ involves only the schedule index $N$, not $p$.
The same observation applies to the prices used below. The threshold
powers, column allowances, profile costs and reserve inverses are
sequences of $N$ alone. Their subpower proofs use the clock and finite
stage purchases; in particular the reserve inverse is bounded by
$4\widehat\beta_N$ and the derivative-shift price by
$\widehat\beta_N$. Their finite products therefore have one subpower
majorant independent of $p$.

\begin{lemma}[Uniform equation sweep and local encoding]
\label{lem:fullmesh-uniform-sweep}
Let $\Pi_{N,p}$ be the uniformly bounded finite projection with kernel
$K_{N,p}$ constructed in
Subsection~\ref{sec:actual-equation-kernel-relay}, using
Subsection~\ref{sec:finite-rank-kernel-transport}.
For $z\in\mathscr C_d$ with $\Pi_{N,p}z_{N,p}=0$, the physical insertion
$O(z)$ has centered root price tending to zero in $I$, at every
sufficiently fine fixed cell depth, in all its finitely many signed and
conjugated orientations. The same holds in translated cells with a fixed
relative interior margin.

Every continuous protected trial whose two prices are uniformly
subpower has, on the same family, a comparison row $x\in\mathscr C_e$
and a discrepancy $O(x)-F$ with such a vanishing root price. In
particular this applies to every finite difference word of a bounded
primary row, with $e$ increased by the word length.
\end{lemma}
\begin{proof}
The columns in \eqref{eq:equation-column-root} are constructed beyond
one $N$-threshold for every $p$, with the same price $E_N$.
All factors in that price are the $N$-only factors just described.
Consequently its subpower bound holds in $I$. For a row in the kernel,
put $v=P_{N,p}^{-1}X_{N,p}^{-1}z$ in the notation of that section;
$|P_{N,p}^{-1}|\le2$. Thus $v$ has a uniform subpower price and the
finite sum of oriented column errors has price
\[
              c_{d,m,k} E_N |v_{N,p}|_1 e^{-s_N/8}.
\]
For each fixed opened pair this is at most $e^{-s_N/16}$ eventually
in $I$. The common-tiling argument there transfers it to any prescribed
coarser cell and relative margin, using a finer fixed opened depth.
Its finite sum and fixed geometric factors can, if needed, be absorbed
by reducing the rate to $e^{-s_N/32}$. The alignment identities and
interiority thresholds depend on $a_N$ and the fixed stencils, not on $p$.

It remains essential to dispose of the comparison insertion itself,
without using an ordered spectator estimate. Since $z\in K_{N,p}$,
the comparison is the campaign insertion of an element of the exact
Maxwell equation ideal: the factors $S_{\theta_N^{-1}}S_{\theta_N}$
cancel exactly. Choose a finite basis of that ideal and one radius
containing all its microscopic equation stencils. In a centered root
test at a fixed physical cell size, the other reflected copies depend
only on links outside this radius once $N$ is large. The smallest
admitted physical period is positive, so the finite nonwrapping floors
are also reached by one $N$-threshold for every $p$. Apply the exact
microscopic integration identity to one copy and the product of all
other copies. This product is bounded and measurable at each finite
cutoff. Its integral is exactly zero, even if the coefficients of the
equation vector are large. The centered root of the comparison is
therefore zero. Linearity covers every vector in the ideal, including
cutoff-dependent vectors; the stencil threshold is independent of those
coefficients. The root triangle inequality now gives the claimed sweep.
No conversion to an ordered spectator limit occurs.

For the second assertion use the terminal construction in
Subsection~\ref{sec:actual-insertion-derivative-relay}. Its price
$E_{N,p}$ is a finite sum of products of the two protected prices and
the same schedule factors. It has a uniform subpower majorant.
The actual descent gives
$|X_{N,p}^{-1}x_{N,p}|_1\le2E_{N,p}$ and root discrepancy
$c_{e,m,k}E_{N,p}e^{-s_N/8}$. The preceding common-tiling argument is
therefore uniform in $p$ here as well. For a word of length $q$, repeat
\eqref{eq:relay-difference-weight} $q$ times. At engineering weight
$d+q$ the powers $a_N^{-q}$ cancel; the stencil is a finite union of
at most $2^q$ fixed translates. The derivative price is at most
$(J_N+1)^q$ times the initial price and the global price at most $2^q$
times it. These are uniformly subpower. This proves the assertion for
all finite words and their orientations without assuming a bound for
the differentiated correlations.
\end{proof}

\begin{lemma}[A uniformly bounded primary and a finite detector]
\label{lem:fullmesh-uniform-detector}
Fix a grade $e$ and a finite direction word $l$. There are finite linear
detectors $L_i$ on its reference coefficient space and $Z>0$, chosen
before $N,p$, such that eventually in $I$
\[
          |\Pi_i v|\le Z|L_i\Pi_i v|\qquad(v\in V_e).
\]
The detectors are averages of correlations with finitely many bounded
primary spectator rows in small neighborhoods of the purchased generic
detector configurations. They have the following two properties:
\begin{enumerate}
\item If $F$ has a uniform paid primary root budget, the assay of
$D_lF$ is bounded eventually in $I$.
\item A continuous trial with a vanishing uniform root price has assay
 tending to zero in $I$.
\end{enumerate}
Every bounded primary row has the required uniform root budget, at
every sufficiently fine fixed centered depth, as do its orientations.
\end{lemma}
\begin{proof}
The last assertion follows directly from
Corollary~\ref{cor:ex-c02-uniform-census}. At a fixed grade and radius
its width $w_k$ is independent of $N,p$ and the label in that ball;
its displayed normalized moment bound gives a finite root budget.
For a requested cell depth choose a finer opened depth and transfer
that budget by the replicated-orbit and maximum-principle argument of
Theorem~\ref{ex-centered-producers}. All thresholds are stencil
interiority thresholds or $N>k$, uniformly in $p$. Finitely many
orientations remain in a bounded ball because their coefficient maps
are fixed linear or conjugate-linear maps.

Write $r=\dim(V_e/\mathcal I_e)$ for the free quotient dimension,
with the campaign identification understood. The exact calibrated
isomorphism taking $\mathcal I_e$ to $K_i$ gives
$\operatorname{rank}\Pi_i=r$ at each finite cutoff; this does not use a
limiting null-space assertion. If $r=0$ the projection vanishes and
there is nothing to prove. Otherwise choose the $r$ independent free
columns and finitely many separating rows used in
Subsection~\ref{sec:actual-rank-purchases}. Their target matrix $G$ is
injective. In supremum norms choose $K$ with $|c|\le K|Gc|$.
If every entry of $Q-G$ is less than $\delta=(2Kr)^{-1}$, then
$|c|\le K|Qc|+\tfrac12|c|$, hence $|c|\le2K|Qc|$.

The purchased comparison in that section holds before extraction:
its corrected root error is $\epsilon_k$, and undoing the finitely
many connected subtractions costs the free-subset moment polynomial
$P_k$ of \eqref{eq:rank-purchased-point-comparison}. At each fixed
comparison level its rescaled lattice configurations are fixed before
$N,p$; let $G_k$ be their free measurement matrix. These configurations
converge to the prescribed generic configurations as $k$ increases.
Continuity of the finitely many free moments on a common separated
neighborhood gives $G_k\to G$ and bounds $P_k$. Choose one sufficiently
deep comparison level $k$ so that, entrywise,
$|G_k-G|<\delta/8$ and $P_k\epsilon_k<\delta/8$.
Before either cutoff limit is taken, the finite-cell inequality applied
to the paid corrected coefficient therefore puts the actual point
measurement matrix within $\delta/4$ of $G$, beyond one mesh threshold
for every period. This is the finite-cutoff estimate used here; its
later passage to the ordered limits is not needed. Its lifts $u_{j,i}$ have
one norm bound $B_k$ for all $N,p$, by
\eqref{eq:rank-fixed-scale-bounded}. Their detector positions have a
fixed interior margin in the comparison cell. Choose small closed
neighborhoods of those positions inside the smallest admitted torus
and the generic set. All coordinate gaps there are at least $4\rho>0$
modulo every period: for the smallest period this follows by a compact
minimum; larger periods cannot create a wrap inside this common
interior box. Shrink the neighborhoods if necessary.

The primary budgets just established give, on these neighborhoods, the
uniform lattice modulus proved in the directional part of
Theorem~\ref{ex-centered-producers}:
\[
 |H_i(x)-H_i(y)|\le 4n A^n\rho^{-1}|x-y|.
\]
Here $n$ is the arity of the detector row, all its labels lie in the
fixed ball containing the lifts, and the centered depth is chosen with
$a_m<\rho$. The proof uses the bounded primary bricks and the
Schatten--Cauchy transfer estimate; it does not use a derivative-field
bound. All gaps and primary budgets here are uniform over every period,
so the displayed proof gives one mesh threshold for this neighborhood.
The exact aligned detector sites represent the same physical centers
once $N$ exceeds their fixed depth.

We specify the averaging operation, also to control derivatives. Choose
$b>0$ small enough and put $h_N=\lceil b/a_N\rceil$. Average a detector
correlation over shifts $0,\ldots,h_N-1$ in the distinguished slot,
once for each direction in $l$. The total physical displacement is
at most $2|l|b$ eventually. Choose $b$ so that this lies inside the
neighborhoods and changes each primary detector entry by less than
$\delta/4$. For the empty word the averaging is the identity.
These averages define $L_i$. The sweep lemma applies to the bounded
rows $(I-\Pi_i)u_{j,i}$: both the rows and their projected parts are
bounded uniformly, and thus belong to $\mathscr C_e$. Planting with
the other bounded detector rows bounds each averaged swept correlation
by a constant times its vanishing root price, uniformly on these
neighborhoods. Consequently the matrix
$Q_i=(L_i\Pi_i u_{j,i})_j$ is within $\delta$ of $G$ eventually in $I$.

It follows that the $r$ vectors $\Pi_i u_{j,i}$ are independent and
therefore span $\operatorname{ran}\Pi_i$. Write
$\Pi_i v=\sum_j c_j\Pi_i u_{j,i}$. If $A$ bounds $\|\Pi_i\|$, then
\[
 |\Pi_i v|\le ArB_k|c|\le 2ArB_kK|L_i\Pi_i v|.
\]
Take $Z\ge1$ at least this last constant. This proves the detector
bound on every vector, not merely on the purchased columns.

For its assay properties, use the same separated neighborhoods and
plant the distinguished trial with the bounded primary spectators in
distinct cells. If its root price is $u_i$, the result is bounded by
$C u_i$, with $C$ independent of $N,p$ and the point in the neighborhood.
Averages have total mass one, so this also proves assertion (2).
For (1), move differences through the expectation by linearity. One
box average telescopes a difference exactly:
\[
 \frac1{h_N}\sum_{j=0}^{h_N-1}
 \frac{H(x+(j+1)e_\alpha)-H(x+je_\alpha)}{a_N}
       =\frac{H(x+h_Ne_\alpha)-H(x)}{h_Na_N}.
\]
The difference and averaging operators commute. Applying this identity
once for each entry of $l$ bounds the averaged derivative by
$(2/b)^{|l|}$ times the undifferentiated correlation bound, since
$h_Na_N\ge b$. That bound follows from the paid root budget of $F$
and the bounded spectator budgets. All translated endpoints remain
in the chosen neighborhoods. This proves (1) without estimating
$D_lF$ pointwise in advance.
\end{proof}

\begin{proposition}[All-period payment for the actual difference trials]
\label{prop:fullmesh-selected-supplier}
The retained public family supplies every coefficient, remainder,
orientation and root-budget hypothesis of
Theorem~\ref{thm:ex-fullmesh-fixedbounds-producer}. Consequently that
theorem's conclusion holds for its actual normalized insertions, with
one mesh threshold valid for every period.
\end{proposition}
\begin{proof}
Let $F$ be a bounded primary row, initially a fixed public label, and
$l$ a finite direction word. At grade $e=d+|l|$, the encoding lemma
gives $x\in\mathscr C_e$ and $E=O(x)-D_lF$ with uniform vanishing
root price. Apply the detector lemma at this grade and word. Its assay
of $D_lF$ is bounded and that of $E$ tends to zero. Thus $|L_ix_i|\le c$
eventually for some $c$ independent of $p$.

Put $y_i=\Pi_i x_i$ and
$z_i=(1+|y_i|)^{-1}(x_i-y_i)$. The first term in this subtraction
belongs to $\mathscr C_e$ by its scalar closure, and the second is
norm-bounded by one. Thus $z\in\mathscr C_e$, and $\Pi_i z_i=0$.
The sweep and assay lemmas give $|L_i z_i|\to0$ in $I$.
Lemma~\ref{lem:ex-fullmesh-cap} now gives $|y_i|\le2Zc+1$ eventually.
Changing $y_i$ to zero at earlier mesh cutoffs makes it a globally
bounded row without affecting any eventual assertion.

This bounded row has a uniform paid primary budget by the detector
lemma. The difference $x-y$ belongs to $\mathscr C_e$ and has exact
projected kernel zero on the tail, so its insertion has a vanishing
uniform root price. The exact identity
\[
               D_lF=O(y)+O(x-y)-E
\]
therefore gives a finite root budget for $D_lF$, at every sufficiently
fine fixed cell depth. Each of its two vanishing prices is eventually
at most one. The root triangle inequality gives the sum of the primary
budget and two, with one common mesh threshold valid for all $p$.
Apply the same argument to the finitely many orientations, or retain
them in the finite lists throughout; take maxima only after fixing the
cell depth. No limit of a moving projection is used.

Finally the periodic difference formula is the definition of these
trials and is exact also at wrapped sites. The planted-cell product
inequality is the Wilson reflection-positive inequality of
Subsection~\ref{ex-centered-estimates}; its chamber interiority follows
from the fixed finite stencils. These verify the remaining hypotheses
of the full-mesh theorem. Its two-scale packing argument now applies
to each finite list of actual insertions and difference directions.
The use of full-mesh estimates has been confined throughout to
uniformly subpower rows and uniform primary budgets; no assertion for
arbitrary ordered subpower rows has been strengthened by changing its
quantifiers.
\end{proof}

\begin{remark}[What this estimate supplies]
The statement is the full-mesh separated-kernel estimate: each finite difference
order and each positive separation have one eventual mesh threshold valid at
all periods. The bounds at fixed period also control compact families of
separated positions and their coordinate ties. The estimate does not itself
give convergence of smeared correlations; that is the separate
same-history convergence estimate. With both inputs, the positional theorem
follows by the interpolation, compactness, and dominated-convergence argument
in Subsection~\ref{sec:ex-limits}.
\end{remark}
